\section{A sharp finite torsion test for simultaneous nonvanishing}
\label{sec:sharp}

The repository controls valuations to within an arbitrarily small error. When
an application needs only nonvanishing, a much smaller explicit bound is
available. The following theorem is developed here from the same
specialization viewpoint. Its proof is independent of RH and GRH. The claim
is a proved mathematical deduction, with sharpness as specified below; a
claim of priority over the literature is not made.

Fix a prime $p$, write $v(p)=1$ and $v(0)=+\infty$, and put
$\mathcal O=\mathcal O_{\mathbb C_p}$. For $r\in\mathbb Q_{>0}$ define
\begin{equation}
 h_*(r)=\min\left\{h\geq1:
       \frac1{p^{h-1}(p-1)}\leq r\right\}.
 \label{eq:hstar}
\end{equation}
In the theorem, different functions are allowed to have different reference
points. All valuations of reference values are finite.

\begin{theorem}[Explicit simultaneous nonvanishing]
\label{thm:sharp-nonvanishing}
Let $s\geq1$ and let $F_1,\ldots,F_q$ be a finite family in
$\mathcal O[[T_1,\ldots,T_s]]$. Suppose that points
$x_i\in\mathbb C_p^s$ satisfy
\[
 v(x_{ia})\geq r>0,\qquad
 B_0:=\sum_{i=1}^q v(F_i(x_i))\leq B<\infty.
\]
Set
\begin{equation}
 H=h_*(r)+\lfloor B\rfloor.
 \label{eq:linear-H}
\end{equation}
Then there are roots of unity $\zeta_1,\ldots,\zeta_s$, each of exact
order $p^{h_a}$ for some $h_*(r)\leq h_a\leq H$, such that
\[
 F_i(\zeta_1-1,\ldots,\zeta_s-1)\ne0\quad(1\leq i\leq q).
\]
Thus $\zeta_a^{p^H}=1$ and $0<v(\zeta_a-1)\leq r$. The depth $H$
depends only on $p,r,B$, not on the number of variables, the number of
functions separately from their aggregate budget, or polynomial degrees.
\end{theorem}

For an empty family the conclusion just means choosing any allowed torsion
tuple. The case $s=0$ reduces to nonzero constants and needs no roots of
unity. These harmless cases will be suppressed in the proof.

\subsection{Weighted Gauss valuation}

For a nonzero integral series $F=\sum_\alpha a_\alpha T^\alpha$ put
\[
 w_r(F)=\min_\alpha\bigl(v(a_\alpha)+r|\alpha|\bigr),
 \qquad w_r(0)=+\infty.
\]
The minimum is attained: coefficients are integral and $r|\alpha|$ tends
to infinity, so only finitely many degrees compete below any fixed finite
value. For integral series one has
\begin{align}
 w_r(F+G)&\geq\min(w_r(F),w_r(G)),\label{eq:gauss-add}\\
 w_r(FG)&=w_r(F)+w_r(G),\label{eq:gauss-mul}\\
 w_r(F)&\leq v(F(x))\qquad\text{if all }v(x_a)\geq r.
 \label{eq:gauss-eval}
\end{align}
The last inequality follows term by term and then by completeness.
For completeness, multiplicativity can be seen by choosing
$u\in\overline{\mathbb Q}_p$ with $v(u)=r$. The series $F(uT)$ has
coefficients tending to zero with total degree. After scaling $F(uT)$ and
$G(uT)$ to make their least coefficient valuations zero, their reductions
are nonzero polynomials over the residue field. Their product is nonzero,
since that polynomial ring is a domain. Hence the product also has least
coefficient valuation zero. Undoing the scaling proves
\eqref{eq:gauss-mul}.

\subsection{Integral division on a Cartesian grid}

The needed division statement does not require that $\mathcal O$ be
Noetherian or discretely valued.

\begin{lemma}[Integral division by a distinguished polynomial]
\label{lem:division}
Suppose $g(T)=T^d+a(T)\in\mathbb Z_p[T]$ is monic, $d\geq1$,
$\deg a<d$, and every coefficient of $a$ is divisible by $p$. Every
$F\in\mathcal O[[T_1,\ldots,T_s]]$ has a decomposition
\[
 F=g(T_1)Q+R,\qquad Q,R\in\mathcal O[[T_1,\ldots,T_s]],
 \qquad\deg_{T_1}R<d.
\]
Successive division in each variable gives
\begin{equation}
 F=\sum_{a=1}^s g(T_a)Q_a+R,
 \qquad\deg_{T_a}R<d\quad(1\leq a\leq s),
 \label{eq:multidivision}
\end{equation}
with all $Q_a$ integral and $R$ a polynomial in all variables.
\end{lemma}

\begin{proof}
Give bounded coefficient arrays the coefficient supremum norm. They form a
complete normed space over $\mathbb C_p$; its integral unit ball is the
set of integral formal series. Let $S_d$ discard the first $d$ powers of
$T_1$ and shift the remaining powers down by $d$. Let $M_a$ multiply by
$a(T_1)$. The operator $K=S_dM_a$ has norm at most $p^{-1}$, since every
coefficient of $a$ has valuation at least one. Therefore
\[
 Q=\sum_{n=0}^{\infty}(-K)^nS_dF
\]
converges in coefficient supremum norm to an integral series. It satisfies
$(1+K)Q=S_dF$. For $R=F-g(T_1)Q$ this identity says $S_dR=0$, so
$\deg_{T_1}R<d$. Division in $T_2$ of this remainder preserves its degree
bound in $T_1$, because the operators act only on $T_2$. Continue through
all variables to obtain \eqref{eq:multidivision}.
\end{proof}

\begin{lemma}[Vanishing on a finite product grid]
\label{lem:grid-ideal}
Under Lemma~\ref{lem:division}, suppose that $g$ has $d$ distinct roots
$A\subset\mathbb C_p$, all of positive valuation. If an integral formal
series $F$ vanishes at every point of $A^s$, then
\[
 F\in\bigl(g(T_1),\ldots,g(T_s)\bigr)
       \mathcal O[[T_1,\ldots,T_s]].
\]
\end{lemma}

\begin{proof}
All integral series in \eqref{eq:multidivision} converge on $A^s$. At a
grid point every $g(T_a)$ vanishes, so the remainder $R$ vanishes there.
It has degree less than $d$ in every variable. Repeated one-variable
interpolation implies $R=0$: fix all but the first coordinate, use $d$
distinct roots, and then repeat for the coefficient polynomials. No
integrality of interpolation denominators is needed; division has already
produced integral quotients before interpolation is used over the field.
\end{proof}

\subsection{Proof and sharpness of the depth bound}

\begin{proof}[Proof of Theorem~\ref{thm:sharp-nonvanishing}]
Write $h_*=h_*(r)$, $L=H-h_*+1=\lfloor B\rfloor+1$, and set
\begin{equation}
 g(T)=\prod_{h=h_*}^{H}\Phi_{p^h}(1+T),
 \label{eq:cyclotomic-grid-poly}
\end{equation}
where $\Phi_n$ is the $n$th cyclotomic polynomial. Modulo $p$,
$\Phi_{p^h}(1+T)=T^{\varphi(p^h)}$, so $g$ is monic with every
nonleading coefficient divisible by $p$. Its roots are precisely
$\zeta-1$ for primitive $p^h$-th roots $\zeta$ with $h_*\leq h\leq H$;
they are distinct and have valuations $\rho_h\leq r$.

For a linear factor $T-\alpha$ one has
$w_r(T-\alpha)=\min(r,v(\alpha))$. By multiplicativity,
\begin{equation}
 w_r(g)=\sum_{h=h_*}^{H}\varphi(p^h)\rho_h=L>B.
 \label{eq:grid-cost}
\end{equation}
Let $F=\prod_iF_i$. Again by multiplicativity and evaluation,
\[
 w_r(F)=\sum_iw_r(F_i)\leq\sum_iv(F_i(x_i))\leq B.
\]
If the desired tuple did not exist, at every point of the grid $A^s$ at
least one $F_i$ would vanish, so $F$ would vanish on the entire grid.
Lemma~\ref{lem:grid-ideal} would give $F=\sum_ag(T_a)Q_a$ with integral
$Q_a$. Since $w_r(Q_a)\geq0$, \eqref{eq:gauss-add} and
\eqref{eq:grid-cost} would imply $w_r(F)\geq L>B$, a contradiction.
\end{proof}

\begin{proposition}[Sharpness for one exponent]
\label{prop:sharpness}
The depth in \eqref{eq:linear-H} cannot be decreased uniformly, even for
one variable and one polynomial, when output roots must satisfy
$0<v(\zeta-1)\leq r$.
\end{proposition}

\begin{proof}
Put $k=\lfloor B\rfloor$. If $k\geq1$, take
\[
 F(T)=\prod_{h=h_*}^{h_*+k-1}\Phi_{p^h}(1+T),\qquad x=0.
\]
Then $v(x)=+\infty\geq r$, while $F(0)=p^k$ has valuation $k\leq B$.
Every permitted torsion point of order dividing $p^{h_*+k-1}$ is a root
of $F$. The next level, $h_*+k$, is therefore necessary. If $k=0$, the
level $h_*$ is already the least one meeting the radial requirement, so
it cannot be reduced either.
\end{proof}

This statement of sharpness concerns $m=1$ and the stated radial
restriction. It is not a claim that every individual family needs this
depth, nor that the first-$m$-powers bound below is optimal.

\subsection{One tuple for the first several powers}

\begin{corollary}[Simultaneous nonvanishing along a finite power orbit]
\label{cor:power-orbit}
Under the assumptions of Theorem~\ref{thm:sharp-nonvanishing}, let
$m\geq1$ and put
\begin{equation}
 H=h_*(r/m)+\lfloor mB\rfloor.
 \label{eq:power-H}
\end{equation}
There is a tuple with $\zeta_a^{p^H}=1$ and
$0<v(\zeta_a-1)\leq r/m$ such that
\[
 F_i(\zeta_1^j-1,\ldots,\zeta_s^j-1)\ne0
 \quad\text{for every }i\text{ and }1\leq j\leq m.
\]
\end{corollary}

\begin{proof}
For each $i,j,a$, choose $y_{ija}$ such that
\[
 (1+y_{ija})^j-1=x_{ia},\qquad v(y_{ija})\geq r/j\geq r/m.
\]
Such a choice exists. If $x_{ia}=0$, use $y_{ija}=0$. Otherwise all roots
of the monic integral polynomial $(1+Y)^j-(1+x_{ia})$ are integral, and
the sum of their valuations is $v(x_{ia})$. At least one root has
valuation at least $v(x_{ia})/j$.

The formal substitutions
\[
 G_{ij}(T)=F_i\bigl((1+T_1)^j-1,\ldots,(1+T_s)^j-1\bigr)
\]
are well defined integral series, since the substituted series have zero
constant term. Their values at $y_{ij}$ equal $F_i(x_i)$; convergence
also follows because all arguments have positive valuation. The aggregate
reference budget for the $mq$ functions $G_{ij}$ is at most $mB$.
Apply Theorem~\ref{thm:sharp-nonvanishing} with radius $r/m$ and budget
$mB$.
\end{proof}

\subsection{Ranks and exact complexes}

\begin{corollary}[Simultaneous rank preservation]
\label{cor:rank}
Let $A_b(T)$ be finitely many matrices over
$\mathcal O[[T_1,\ldots,T_s]]$. For each $b$ select an $r_b$-minor
$D_b(T)$ and an interior point $x_b$ with
$v(D_b(x_b))<\infty$. Put $B=\sum_bv(D_b(x_b))$, and assume all
coordinates of all $x_b$ have valuation at least $r>0$.
With $H$ as in \eqref{eq:power-H}, one tuple of torsion parameters gives
\[
 \operatorname{rank}_{\mathbb C_p} A_b(\zeta^j-1)\geq r_b
 \qquad(b,\ 1\leq j\leq m).
\]
If $r_b$ is the rank over the fraction field of the formal-series domain,
all these inequalities are equalities.
\end{corollary}

\begin{proof}
Apply Corollary~\ref{cor:power-orbit} to the selected determinants. For
the equality assertion, all larger minors are identically zero formal
series and hence vanish at every convergent specialization.
\end{proof}

For a finite complex of finite free modules, apply the corollary to one
generic-rank minor of each differential. Every resulting complex over
$\mathbb C_p$ has its generic ranks, so its cohomology dimensions are the
generic dimensions. In particular, a generically exact complex becomes
exact at every one of the first $m$ powers of the selected torsion tuple.
This conclusion concerns cohomology after passing to the field; it does
not assert that integral torsion in cohomology vanishes.

\subsection{Essential limitations and an exact example}

\paragraph{Integral coefficients are essential.}
The analytic function $\log(1+T)=\sum_{n\geq1}(-1)^{n+1}T^n/n$
vanishes at every $p$-power torsion point $\zeta-1$, but is nonzero.
Its coefficients do not have a uniform lower valuation bound. It therefore
does not satisfy the theorem. Conversely, any family with a known common
lower coefficient bound can be scaled to an integral family, but that
scaling changes the reference budget and must be included.

\paragraph{Existence is weaker than near preservation of valuations.}
The linear depth \eqref{eq:power-H} guarantees that the selected values
are finite, not that they are within $\varepsilon$ of their reference
valuations. The stronger conclusion is addressed separately in
Section~\ref{sec:series-specialization}. Neither theorem bounds the norm of
the conductor of a global arithmetic character without an additional
construction relating these torsion parameters to that character.

\paragraph{Exact finite illustration.}
Take $p=3$, $r=1$, $m=3$, and
\[
 A(X,Y)=\begin{pmatrix}1&X\\Y&XY+X\end{pmatrix},\qquad
 B(X,Y)=\begin{pmatrix}X-Y&0\\0&1\end{pmatrix}.
\]
Their determinants are $X$ and $X-Y$. Add the function
$F_3(X,Y)=(1+X)^3-1$. Reference points $(3,3)$, $(3,6)$, and $(3,3)$
give valuations $1,1,2$, so $B=4$. The guaranteed depth is
$h_*(1/3)+12=14$. The explicit tuple
$(\zeta_{27},\zeta_{27}^2)$ already works at depth $3$:
\begin{center}
\begin{tabular}{cccc}
\toprule
Power $j$ & $v(\det A)$ & $v(\det B)$ & $v(F_3)$\\
\midrule
1 & $1/18$ & $1/18$ & $1/6$\\
2 & $1/18$ & $1/18$ & $1/6$\\
3 & $1/6$ & $1/6$ & $1/2$\\
\bottomrule
\end{tabular}
\end{center}
Every entry is finite, so both matrices have rank two at all three powers.
The accompanying script checks these equalities by exact cyclotomic
resultants, and verifies eighteen finite sharpness examples. These examples
illustrate the theorems; the proofs above establish the universal assertions.
